\subsection{CANONICAL EXTENSIONS OF MAPPINGS}

Let \(S = (c_1, c_2, \ldots, c_m)\) be an echelon construction scheme on \(n\) terms. Let \(E_1, \ldots, E_n, E_1', \ldots, E_n'\) be sets (terms in \(\mathcal{G}\) ) and let \(f_1, \ldots, f_n\) be terms in \(\mathcal{G}\) such that the relations `` \(f_i\) is a mapping of \(E_i\) into \(E_i'\) '' are theorems in \(\mathcal{G}\) for \(1 \leqslant i \leqslant n\) . Let \(A_1, \ldots, A_m\) (resp. \(A_1', \ldots, A_m'\) ) be the echelon construction of scheme S on \(E_1, \ldots, E_n\) (resp. \(E_1', \ldots, E_n'\) ). We define step by step a sequence of \(m\) terms \(g_1, \ldots, g_m\) such that \(g_i\) is a mapping of \(A_i\) into \(A_i'\) (for \(1 \leqslant i \leqslant m\) ) by the following conditions:

\begin{enumerate}
\def\labelenumi{(\alph{enumi})}
\item
  If \(c_{i} = (0, b_{i})\) , so that \(A_{i} = E_{b_{i}}\) and \(A_{i}' = E_{b_{i}}'\) , then \(g_{i}\) is the mapping \(f_{b_{i}}\) .
\item
  If \(c_{i} = (a_{i}, 0)\) , so that \(\mathbf{A}_{i} = \mathfrak{P}(\mathbf{A}_{a_{i}})\) and \(\mathbf{A}_{i}' = \mathfrak{P}(\mathbf{A}_{a_{i}}')\) , then \(g_{i}\) is the canonical extension \(\hat{g}_{a_i}\) of \(g_{a_i}\) to sets of subsets (Chapter II, § 5, no. 1).
\item
  If \(c_{i} = (a_{i}, b_{i})\) , where \(a_{i} \neq 0\) and \(b_{i} \neq 0\) , so that
\end{enumerate}

\[
\mathrm{A} _ {i} = \mathrm{A} _ {a _ {i}} \times \mathrm{A} _ {b _ {i}} \quad \text { and } \quad \mathrm{A} _ {i} ^ {\prime} = \mathrm{A} _ {a _ {i}} ^ {\prime} \times \mathrm{A} _ {b _ {i}} ^ {\prime},
\]

then \(g_{i}\) is the canonical extension \(g_{a_i} \times g_{b_i}\) of \(g_{a_i}\) and \(g_{b_i}\) to \(A_{a_i} \times A_{b_i}\) (Chapter II, § 3, no. 9).

The last term \(g_{m}\) of this sequence is called the canonical extension, with scheme S, of the mappings \(f_{1}, \ldots, f_{n}\) , and will be denoted by \(\langle f_{1}, \ldots, f_{n} \rangle^{S}\) . The following criteria can be verified step by step:

CST1. If \(f_{i}\) is a mapping of \(\mathbf{E}_i\) into \(\mathbf{E}_i'\) , and if \(f_{i}'\) is a mapping of \(\mathbf{E}_i'\) into \(\mathbf{E}_i''\) ( \(1 \leqslant i \leqslant n\) ), then for every echelon construction scheme S on \(n\) terms we have

\[
\langle f _ {1} ^ {\prime} \circ f _ {1}, f _ {2} ^ {\prime} \circ f _ {2}, \dots , f _ {n} ^ {\prime} \circ f _ {n} \rangle^ {\mathrm{S}} = \langle f _ {1} ^ {\prime}, f _ {2} ^ {\prime}, \dots , f _ {n} ^ {\prime} \rangle^ {\mathrm{S}} \circ \langle f _ {1}, f _ {2}, \dots , f _ {n} \rangle^ {\mathrm{S}}.
\]

CST2. If \(f_{i}\) is injective (resp. surjective) for \(1 \leqslant i \leqslant n\) , then \(\langle f_{1}, \ldots, f_{n} \rangle^{S}\) is injective (resp. surjective).

This criterion follows from the corresponding properties of the extension \(\hat{g}\) (Chapter II, § 5, no. 1, Proposition 1) and the extension \(g \times h\) (Chapter II, § 3, no. 9).

CST3. If \(f_{i}\) is a bijection of \(\mathbf{E}_i\) onto \(\mathbf{E}_i'\) , and if \(f_{i}^{-1}\) is the inverse bijection (*), then \(\langle f_1, \ldots, f_n \rangle^{S}\) is a bijection and \(\langle f_1^{-1}, \ldots, f_n^{-1} \rangle^{S}\) is its inverse; in other words,

\[
(\langle f _ {1}, \dots , f _ {n} \rangle^ {\mathrm{S}}) ^ {- 1} = \langle f _ {1} ^ {- 1}, \dots , f _ {n} ^ {- 1} \rangle^ {\mathrm{S}}.
\]

This follows immediately from CST1 and CST2.

